% =====================================================================
% 4. Modelagem do Grafo
% Responsável: Pessoa 1 (Maria Eduarda) — ESCRITA
%
% Baseada em docs/modelo_lpg.md e schema/schema.cypher. Se o modelo
% mudar, atualizar as tabelas e a figura.
% =====================================================================

\section{Modelagem do Grafo}
\label{sec:modelagem}

O conhecimento extraído é armazenado no Neo4j segundo o modelo de grafo de propriedades rotuladas (\emph{Labeled Property Graph}, LPG) \cite{angles2018pg,robinson2015graphdb}. Nesse modelo, nós e relações têm um tipo (rótulo) e podem carregar propriedades na forma chave--valor, o que permite representar tanto as entidades do domínio quanto informações sobre a origem de cada fato.

\subsection{Duas camadas: conhecimento e procedência}

O grafo foi dividido em duas camadas, ilustradas na Figura~\ref{fig:modelo}. A \emph{camada de conhecimento} contém as entidades do domínio e as relações entre elas; é sobre ela que ocorre a travessia multi-hop. A \emph{camada de procedência} contém os documentos da Wikipédia (\rot{Documento}) e os trechos em que foram divididos (\rot{Chunk}). As duas camadas se ligam pela relação \rel{MENCIONA}, de um trecho para cada entidade citada nele, e pela relação \rel{SOBRE}, de um documento para a entidade de que ele trata.

Essa separação atende diretamente ao objetivo de reduzir alucinações: qualquer entidade ou relação usada para responder a uma pergunta pode ser rastreada até o trecho original de onde foi extraída, e esse trecho pode ser mostrado ao usuário como evidência.

\begin{figure}[ht]
\centering
\resizebox{\linewidth}{!}{%
\begin{tikzpicture}[
  ent/.style={draw, rounded corners=3pt, minimum width=2.2cm, minimum height=0.7cm, font=\footnotesize},
  prov/.style={ent, fill=gray!12},
  lab/.style={font=\scriptsize, inner sep=1.5pt},
  >={Stealth[length=1.6mm]}]
  % camada de procedência (coluna da esquerda)
  \node[prov] (doc) at (0,0)  {Documento};
  \node[prov] (chk) at (0,-2) {Chunk};
  % camada de conhecimento (grade 2 x 3)
  \node[ent] (pes) at (5.0,0)  {Pessoa};
  \node[ent] (org) at (10.3,0)  {Organizacao};
  \node[ent] (lin) at (5.0,-2) {Linguagem};
  \node[ent] (tec) at (10.3,-2) {Tecnologia};
  \node[ent] (hw)  at (10.3,-4) {Hardware};
  % relações de procedência e ligação entre camadas
  \draw[->] (doc) -- node[lab, left] (tc) {TEM\_CHUNK} (chk);
  \draw[->, dashed] (doc) -- node[lab, above] {SOBRE} (pes);
  \draw[->, dashed] (chk) -- node[lab, above] {MENCIONA} (lin);
  % relações da camada de conhecimento
  \draw[->] (pes) -- node[lab, above] {TRABALHOU\_EM} (org);
  \draw[->] (pes) -- node[lab, left] {CRIOU} (lin);
  \draw[->] (org) -- node[lab, right] (dv) {DESENVOLVEU} (tec);
  \draw[->] (tec) -- node[lab, above] {ESCRITO\_EM} (lin);
  \draw[->] (tec) -- node[lab, right] (ex) {EXECUTA\_EM} (hw);
  % caixas das camadas
  \begin{scope}[on background layer]
    \node[draw, dotted, inner sep=8pt, fit=(doc)(chk)(tc),
          label={[font=\scriptsize]above:Procedência}] {};
    \node[draw, dotted, inner sep=8pt, fit=(pes)(org)(hw)(dv)(ex),
          label={[font=\scriptsize]above:Conhecimento}] {};
  \end{scope}
\end{tikzpicture}%
}
\caption{Visão simplificada do modelo. Linhas tracejadas ligam as duas camadas; \rel{MENCIONA} e \rel{SOBRE} podem apontar para qualquer entidade. Os rótulos \rot{Premio} e \rot{Local} e parte das relações foram omitidos para simplificar.}
\label{fig:modelo}
\end{figure}

\subsection{Entidades e relações}

A camada de conhecimento tem sete tipos de entidade, descritos na Tabela~\ref{tab:entidades}. Além do rótulo específico, todo nó de entidade recebe o rótulo genérico \rot{Entidade}. Isso permite manter um único índice textual e um único índice vetorial para todas as entidades, de modo que a busca não precisa saber de antemão o tipo do que procura. Todas as entidades compartilham as propriedades \texttt{id}, \texttt{nome}, \texttt{aliases} (formas alternativas do nome), \texttt{descricao}, \texttt{wikidata\_id} e \texttt{embedding}.

\begin{table}[ht]
\centering
\caption{Tipos de entidade da camada de conhecimento.}
\label{tab:entidades}
\small
\begin{tabularx}{\textwidth}{lXl}
\toprule
\textbf{Rótulo} & \textbf{Representa} & \textbf{Propriedades específicas} \\
\midrule
\rot{Pessoa}      & Pesquisadores, engenheiros e fundadores & anos de nascimento e morte \\
\rot{Organizacao} & Empresas, universidades, laboratórios, agências e fundações & tipo, ano de fundação \\
\rot{Linguagem}   & Linguagens de programação & ano, paradigmas \\
\rot{Hardware}    & Computadores, processadores e periféricos & ano, tipo \\
\rot{Tecnologia}  & Sistemas, redes, protocolos, invenções e conceitos & ano, tipo \\
\rot{Premio}      & Prêmios e distinções & -- \\
\rot{Local}       & Países e cidades & tipo \\
\bottomrule
\end{tabularx}
\end{table}

Os rótulos \rot{Premio} e \rot{Local} não correspondem a páginas do corpus, mas foram incluídos porque aparecem com frequência nos textos e geram consultas multi-hop naturais, como ``quais ganhadores do Prêmio Turing trabalharam no Bell Labs?''.

Foram definidos 14 tipos de relação, cada um restrito a tipos específicos de origem e destino (Tabela~\ref{tab:relacoes}). Essas restrições funcionam como um esquema semântico: uma relação \rel{CRIOU}, por exemplo, só pode ligar uma pessoa a uma linguagem, um hardware ou uma tecnologia. Extrações que violam essas regras são rejeitadas antes de chegar ao banco, o que evita que erros do modelo de linguagem se tornem fatos no grafo.

\begin{table}[ht]
\centering
\caption{Tipos de relação e restrições de origem e destino.}
\label{tab:relacoes}
\small
\begin{tabularx}{\textwidth}{llX}
\toprule
\textbf{Relação} & \textbf{Origem} & \textbf{Destino} \\
\midrule
\rel{TRABALHOU\_EM} & Pessoa & Organizacao \\
\rel{ESTUDOU\_EM} & Pessoa & Organizacao \\
\rel{FUNDOU} & Pessoa & Organizacao \\
\rel{CRIOU} & Pessoa & Linguagem, Hardware, Tecnologia \\
\rel{DESENVOLVEU} & Organizacao & Linguagem, Hardware, Tecnologia \\
\rel{INFLUENCIADA\_POR} & Linguagem & Linguagem \\
\rel{ESCRITO\_EM} & Tecnologia & Linguagem \\
\rel{EXECUTA\_EM} & Tecnologia, Linguagem & Hardware \\
\rel{UTILIZA} & Hardware & Hardware \\
\rel{SUCESSOR\_DE} & qualquer & mesmo tipo da origem \\
\rel{RECEBEU} & Pessoa & Premio \\
\rel{COLABOROU\_COM} & Pessoa & Pessoa \\
\rel{NASCEU\_EM} & Pessoa & Local \\
\rel{SEDIADA\_EM} & Organizacao & Local \\
\bottomrule
\end{tabularx}
\end{table}

Toda relação extraída guarda três propriedades de procedência: \texttt{chunk\_ids}, a lista de trechos em que foi encontrada; \texttt{evidencia}, a frase do texto que a comprova; e \texttt{confianca}, um valor entre 0 e 1 atribuído na extração. Quando a mesma relação é encontrada em mais de um trecho, ela não é duplicada: o novo trecho é acrescentado à lista existente.

\subsection{Identificadores, restrições e índices}

Cada entidade recebe um identificador canônico formado pelo tipo e pelo nome normalizado, sem acentos e em minúsculas (por exemplo, \texttt{pessoa:alan\_turing} e \texttt{linguagem:cpp} para C++). Como todos os documentos usam a mesma regra, uma entidade citada em várias páginas converge para um único nó já na ingestão. Os casos que a regra não resolve, como ``Turing'' e ``Alan Turing'', são tratados na etapa de resolução de entidades (Seção~\ref{sec:ingestao}).

A integridade do grafo é garantida por 11 restrições de unicidade: sobre o \texttt{id} de cada um dos sete rótulos específicos, do rótulo \rot{Entidade}, de \rot{Documento} e de \rot{Chunk}, e sobre a URL de \rot{Documento}. Para as consultas, foram criados três tipos de índice, resumidos na Tabela~\ref{tab:indices}.

\begin{table}[ht]
\centering
\caption{Índices criados no Neo4j.}
\label{tab:indices}
\small
\begin{tabularx}{\textwidth}{lXX}
\toprule
\textbf{Tipo} & \textbf{Onde} & \textbf{Uso} \\
\midrule
Range (B-tree) & nome das entidades; categoria e título dos documentos; ano & Filtros e buscas exatas em consultas Cypher \\
Fulltext & nome, \texttt{aliases} e descrição das entidades; texto dos trechos & Busca por palavras-chave e por formas alternativas de nomes \\
Vetorial (cosseno) & \texttt{embedding} das entidades e dos trechos & Busca semântica a partir da pergunta do usuário \\
\bottomrule
\end{tabularx}
\end{table}

Os índices vetoriais têm 768 dimensões, correspondentes ao modelo de \emph{embedding} \texttt{multilingual-e5-base} \cite{wang2024me5}. O modelo foi escolhido por ser gratuito, executável localmente, treinado para vários idiomas, incluindo o português, e por aceitar entradas de até 512 \emph{tokens}, o que comporta os trechos do corpus sem truncamento.

\subsection{Verificação do esquema}

O esquema é criado por um script Cypher idempotente, que pode ser executado várias vezes sem duplicar estruturas. Um segundo script verifica o resultado: confere a existência das 11 restrições e dos índices (17 do tipo range, incluindo os criados automaticamente pelas restrições, 2 fulltext e 2 vetoriais, todos no estado \texttt{ONLINE}), insere nós de teste, executa uma consulta em cada índice textual e vetorial e remove os nós de teste ao final. Na versão 5.26 do Neo4j utilizada no projeto, todos os testes passaram, inclusive a busca textual por uma forma alternativa do nome armazenada na lista \texttt{aliases}.
